% Source: book pp. 2--11 (PDF 16--25); solutions pp. 3--8.
\section{\texorpdfstring{$\R^n$ and $\C^n$}{R^n and C^n}}

\subsection{Complex Numbers}

You should already be familiar with basic properties of the set $\R$ of real
numbers. Complex numbers were invented so that we can take square roots of
negative numbers. The idea is to assume we have a square root of $-1$, denoted
by $i$, that obeys the usual rules of arithmetic. Here are the formal
definitions.

\begin{definition}[Complex numbers, $\C$]\label{ax:1.1}\leavevmode
\begin{itemize}
\item A \emph{complex number} is an ordered pair $(a,b)$, where
  $a,b\in\R$, but we will write this as $a+bi$.
\item The set of all complex numbers is denoted by $\C$:
  \[ \C = \{a+bi : a,b\in\R\}. \]
\item \emph{Addition} and \emph{multiplication} on $\C$ are defined by
  \begin{align*}
    (a+bi)+(c+di) &= (a+c)+(b+d)i,\\
    (a+bi)(c+di)  &= (ac-bd)+(ad+bc)i;
  \end{align*}
  here $a,b,c,d\in\R$.
\end{itemize}
\end{definition}

If $a\in\R$, we identify $a+0i$ with the real number $a$. Thus we think of
$\R$ as a subset of $\C$. We usually write $0+bi$ as just $bi$, and we usually
write $0+1i$ as just $i$.

\marginnote{The symbol $i$ was first used to denote $\sqrt{-1}$ by Leonhard
Euler in 1777.}%
To motivate the definition of complex multiplication given above, pretend that
we knew that $i^2=-1$ and then use the usual rules of arithmetic to derive the
formula above for the product of two complex numbers. Then use that formula to
verify that we indeed have
\[ i^2 = -1. \]

Do not memorize the formula for the product of two complex numbers---you can
always rederive it by recalling that $i^2=-1$ and then using the usual rules of
arithmetic (as given by \ref{ax:1.3}). The next example illustrates this
procedure.

\begin{example}[Complex arithmetic]\label{ax:1.2}
The product $(2+3i)(4+5i)$ can be evaluated by applying the distributive and
commutative properties from \ref{ax:1.3}:
\begin{align*}
(2+3i)(4+5i) &= 2\cdot(4+5i) + (3i)(4+5i)\\
             &= 2\cdot4 + 2\cdot5i + 3i\cdot4 + (3i)(5i)\\
             &= 8 + 10i + 12i - 15\\
             &= -7 + 22i.
\end{align*}
\end{example}

Our first result states that complex addition and complex multiplication have
the familiar properties that we expect.

\begin{result}[Properties of complex arithmetic]\label{ax:1.3}\leavevmode
\begin{axprops}
\item[commutativity] $\alpha+\beta=\beta+\alpha$ and $\alpha\beta=\beta\alpha$
  for all $\alpha,\beta\in\C$.
\item[associativity] $(\alpha+\beta)+\lambda=\alpha+(\beta+\lambda)$ and
  $(\alpha\beta)\lambda=\alpha(\beta\lambda)$ for all $\alpha,\beta,\lambda\in\C$.
\item[identities] $\lambda+0=\lambda$ and $\lambda1=\lambda$ for all $\lambda\in\C$.
\item[additive inverse] For every $\alpha\in\C$, there exists a unique
  $\beta\in\C$ such that $\alpha+\beta=0$.
\item[multiplicative inverse] For every $\alpha\in\C$ with $\alpha\ne0$, there
  exists a unique $\beta\in\C$ such that $\alpha\beta=1$.
\item[distributive property] $\lambda(\alpha+\beta)=\lambda\alpha+\lambda\beta$
  for all $\lambda,\alpha,\beta\in\C$.
\end{axprops}
\end{result}

The properties above are proved using the familiar properties of real numbers
and the definitions of complex addition and multiplication. The next example
shows how commutativity of complex multiplication is proved. Proofs of the
other properties above are left as exercises.

\begin{example}[Commutativity of complex multiplication]\label{ax:1.4}
To show that $\alpha\beta=\beta\alpha$ for all $\alpha,\beta\in\C$, suppose
\[ \alpha = a+bi \quad\text{and}\quad \beta = c+di, \]
where $a,b,c,d\in\R$. Then the definition of multiplication of complex numbers
shows that
\begin{align*}
\alpha\beta &= (a+bi)(c+di)\\
            &= (ac-bd)+(ad+bc)i
\end{align*}
and
\begin{align*}
\beta\alpha &= (c+di)(a+bi)\\
            &= (ca-db)+(cb+da)i.
\end{align*}
The equations above and the commutativity of multiplication and addition of
real numbers show that $\alpha\beta=\beta\alpha$.
\end{example}

Next, we define the additive and multiplicative inverses of complex numbers,
and then use those inverses to define subtraction and division operations with
complex numbers.

\begin{definition}[$-\alpha$, subtraction, $1/\alpha$, division]\label{ax:1.5}\leavevmode
Suppose $\alpha,\beta\in\C$.
\begin{itemize}
\item Let $-\alpha$ denote the additive inverse of $\alpha$. Thus $-\alpha$ is
  the unique complex number such that
  \[ \alpha+(-\alpha)=0. \]
\item \emph{Subtraction} on $\C$ is defined by
  \[ \beta-\alpha = \beta+(-\alpha). \]
\item For $\alpha\ne0$, let $1/\alpha$ and $1/\alpha$ denote the
  multiplicative inverse of $\alpha$. Thus $1/\alpha$ is the unique complex
  number such that
  \[ \alpha(1/\alpha)=1. \]
\item For $\alpha\ne0$, \emph{division} by $\alpha$ is defined by
  \[ \beta/\alpha = \beta(1/\alpha). \]
\end{itemize}
\end{definition}

So that we can conveniently make definitions and prove theorems that apply to
both real and complex numbers, we adopt the following notation.

\begin{notation}[$\F$]\label{ax:1.6}
Throughout this book, $\F$ stands for either $\R$ or $\C$.
\end{notation}

\marginnote{The letter $\F$ is used because $\R$ and $\C$ are examples of what
are called \emph{fields}.}%
Thus if we prove a theorem involving $\F$, we will know that it holds when
$\F$ is replaced with $\R$ and when $\F$ is replaced with $\C$.

Elements of $\F$ are called \emph{scalars}. The word ``scalar'' (which is just
a fancy word for ``number'') is often used when we want to emphasize that an
object is a number, as opposed to a vector (vectors will be defined soon).

For $\alpha\in\F$ and $m$ a positive integer, we define $\alpha^m$ to denote
the product of $\alpha$ with itself $m$ times:
\[ \alpha^m = \underbrace{\alpha\cdots\alpha}_{m\text{ times}}. \]
This definition implies that
\[ (\alpha^m)^n = \alpha^{mn} \quad\text{and}\quad (\alpha\beta)^m = \alpha^m\beta^m \]
for all $\alpha,\beta\in\F$ and all positive integers $m,n$.

\subsection{Lists}

Before defining $\R^n$ and $\C^n$, we look at two important examples.

\begin{example}[$\R^2$ and $\R^3$]\label{ax:1.7}\leavevmode
\begin{itemize}
\item The set $\R^2$, which you can think of as a plane, is the set of all
  ordered pairs of real numbers:
  \[ \R^2 = \{(x,y) : x,y\in\R\}. \]
\item The set $\R^3$, which you can think of as ordinary space, is the set of
  all ordered triples of real numbers:
  \[ \R^3 = \{(x,y,z) : x,y,z\in\R\}. \]
\end{itemize}
\end{example}

To generalize $\R^2$ and $\R^3$ to higher dimensions, we first need to discuss
the concept of lists.

\begin{definition}[List, length]\label{ax:1.8}\leavevmode
\begin{itemize}
\item Suppose $n$ is a nonnegative integer. A \emph{list} of \emph{length} $n$
  is an ordered collection of $n$ elements (which might be numbers, other
  lists, or more abstract objects).
\item Two lists are equal if and only if they have the same length and the
  same elements in the same order.
\end{itemize}
\end{definition}

\marginnote{Many mathematicians call a list of length $n$ an
\emph{$n$-tuple}.}%
Lists are often written as elements separated by commas and surrounded by
parentheses. Thus a list of length two is an ordered pair that might be written
as $(a,b)$. A list of length three is an ordered triple that might be written
as $(x,y,z)$. A list of length $n$ might look like this:
\[ (z_1,\dots,z_n). \]

Sometimes we will use the word \emph{list} without specifying its length.
Remember, however, that by definition each list has a finite length that is a
nonnegative integer. Thus an object that looks like $(x_1,x_2,\dots)$, which
might be said to have infinite length, is not a list.

A list of length $0$ looks like this: $(\,)$. We consider such an object to be
a list so that some of our theorems will not have trivial exceptions.

Lists differ from finite sets in two ways: in lists, order matters and
repetitions have meaning; in sets, order and repetitions are irrelevant.

\begin{example}[Lists versus sets]\label{ax:1.9}\leavevmode
\begin{itemize}
\item The lists $(3,5)$ and $(5,3)$ are not equal, but the sets $\{3,5\}$ and
  $\{5,3\}$ are equal.
\item The lists $(4,4)$ and $(4,4,4)$ are not equal (they do not have the same
  length), although the sets $\{4,4\}$ and $\{4,4,4\}$ both equal the set
  $\{4\}$.
\end{itemize}
\end{example}

\subsection{\texorpdfstring{$\F^n$}{F^n}}

To define the higher-dimensional analogues of $\R^2$ and $\R^3$, we will simply
replace $\R$ with $\F$ (which equals $\R$ or $\C$) and replace the $2$ or $3$
with an arbitrary positive integer.

\begin{notation}[$n$]\label{ax:1.10}
Fix a positive integer $n$ for the rest of this chapter.
\end{notation}

\begin{definition}[$\F^n$, coordinate]\label{ax:1.11}
$\F^n$ is the set of all lists of length $n$ of elements of $\F$:
\[ \F^n = \{(x_1,\dots,x_n) : x_k\in\F \text{ for } k=1,\dots,n\}. \]
For $(x_1,\dots,x_n)\in\F^n$ and $k\in\{1,\dots,n\}$, we say that $x_k$ is the
$k^{\text{th}}$ \emph{coordinate} of $(x_1,\dots,x_n)$.
\end{definition}

If $\F=\R$ and $n$ equals $2$ or $3$, then the definition above of $\F^n$
agrees with our previous notions of $\R^2$ and $\R^3$.

\begin{example}[$\C^4$]\label{ax:1.12}
$\C^4$ is the set of all lists of four complex numbers:
\[ \C^4 = \{(z_1,z_2,z_3,z_4) : z_1,z_2,z_3,z_4\in\C\}. \]
\end{example}

\marginnote[-9mm]{Read \emph{Flatland: A Romance of Many Dimensions}, by Edwin A.
Abbott, for an amusing account of how $\R^3$ would be perceived by creatures
living in $\R^2$. This novel, published in 1884, may help you imagine a
physical space of four or more dimensions.}%
If $n\ge4$, we cannot visualize $\R^n$ as a physical object. Similarly, $\C^1$
can be thought of as a plane, but for $n\ge2$, the human brain cannot provide a
full image of $\C^n$. However, even if $n$ is large, we can perform algebraic
manipulations in $\F^n$ as easily as in $\R^2$ or $\R^3$. For example, addition
in $\F^n$ is defined as follows.

\begin{definition}[Addition in $\F^n$]\label{ax:1.13}
\emph{Addition} in $\F^n$ is defined by adding corresponding coordinates:
\[ (x_1,\dots,x_n)+(y_1,\dots,y_n) = (x_1+y_1,\dots,x_n+y_n). \]
\end{definition}

Often the mathematics of $\F^n$ becomes cleaner if we use a single letter to
denote a list of $n$ numbers, without explicitly writing the coordinates. For
example, the next result is stated with $x$ and $y$ in $\F^n$ even though the
proof requires the more cumbersome notation of $(x_1,\dots,x_n)$ and
$(y_1,\dots,y_n)$.

\begin{result}[Commutativity of addition in $\F^n$]\label{ax:1.14}
If $x,y\in\F^n$, then $x+y=y+x$.
\end{result}

\begin{proof}
Suppose $x=(x_1,\dots,x_n)\in\F^n$ and $y=(y_1,\dots,y_n)\in\F^n$. Then
\begin{align*}
x+y &= (x_1,\dots,x_n)+(y_1,\dots,y_n)\\
    &= (x_1+y_1,\dots,x_n+y_n)\\
    &= (y_1+x_1,\dots,y_n+x_n)\\
    &= (y_1,\dots,y_n)+(x_1,\dots,x_n)\\
    &= y+x,
\end{align*}
where the second and fourth equalities above hold because of the definition of
addition in $\F^n$ and the third equality holds because of the usual
commutativity of addition in $\F$.
\end{proof}

\marginnote{The symbol $\blacksquare$ means ``end of proof''.}%
If a single letter is used to denote an element of $\F^n$, then the same letter
with appropriate subscripts is often used when coordinates must be displayed.
For example, if $x\in\F^n$, then letting $x$ equal $(x_1,\dots,x_n)$ is good
notation, as shown in the proof above. Even better, work with just $x$ and
avoid explicit coordinates when possible.

\begin{notation}[$0$]\label{ax:1.15}
Let $0$ denote the list of length $n$ whose coordinates are all $0$:
\[ 0 = (0,\dots,0). \]
\end{notation}

Here we are using the symbol $0$ in two different ways---on the left side of
the equation above, the symbol $0$ denotes a list of length $n$, which is an
element of $\F^n$, whereas on the right side, each $0$ denotes a number. This
potentially confusing practice actually causes no problems because the context
should always make clear which $0$ is intended.

\begin{example}[Context determines which $0$ is intended]\label{ax:1.16}
Consider the statement that $0$ is an additive identity for $\F^n$:
\[ x+0=x \quad\text{for all } x\in\F^n. \]
Here the $0$ above is the list defined in \ref{ax:1.15}, not the number $0$,
because we have not defined the sum of an element of $\F^n$ (namely, $x$) and
the number $0$.
\end{example}

\marginfig[-2mm]{\begin{tikzpicture}[scale=.5]
  \draw[axis] (-.4,0)--(3.6,0); \draw[axis] (0,-.4)--(0,2.8);
  \draw[vec] (0,0)--(2.6,1.9);
  \fill (2.6,1.9) circle (1.4pt) node[above left=-1pt,font=\scriptsize] {$(a,b)$};
\end{tikzpicture}}{Elements of $\R^2$ can be thought of as points or as vectors.}%
A picture can aid our intuition. We will draw pictures in $\R^2$ because we can
sketch this space on two-dimensional surfaces such as paper and computer
screens. A typical element of $\R^2$ is a point $v=(a,b)$. Sometimes we think
of $v$ not as a point but as an arrow starting at the origin and ending at
$(a,b)$, as shown here. When we think of an element of $\R^2$ as an arrow, we
refer to it as a \emph{vector}.

\marginfig[1mm]{\begin{tikzpicture}[scale=.5]
  \draw[vec] (0,0)--(2.4,1.3); \draw[vec] (.9,1.6)--(3.3,2.9);
\end{tikzpicture}}{A vector.}%
When we think of vectors in $\R^2$ as arrows, we can move an arrow parallel to
itself (not changing its length or direction) and still think of it as the same
vector. With that viewpoint, you will often gain better understanding by
dispensing with the coordinate axes and the explicit coordinates and just
thinking of the vector, as shown in the figure here. The two arrows shown here
have the same length and same direction, so we think of them as the same
vector.

\marginnote{Mathematical models of the economy can have thousands of
variables, say $x_1,\dots,x_{5000}$, which means that we must work in
$\R^{5000}$. Such a space cannot be dealt with geometrically. However, the
algebraic approach works well. Thus our subject is called \emph{linear
algebra}.}%
Whenever we use pictures in $\R^2$ or use the somewhat vague language of points
and vectors, remember that these are just aids to our understanding, not
substitutes for the actual mathematics that we will develop. Although we cannot
draw good pictures in high-dimensional spaces, the elements of these spaces are
as rigorously defined as elements of $\R^2$.

For example, $(2,-3,17,\pi,\sqrt2)$ is an element of $\R^5$, and we may
casually refer to it as a point in $\R^5$ or a vector in $\R^5$ without
worrying about whether the geometry of $\R^5$ has any physical meaning.

Recall that we defined the sum of two elements of $\F^n$ to be the element of
$\F^n$ obtained by adding corresponding coordinates; see \ref{ax:1.13}. As we
will now see, addition has a simple geometric interpretation in the special
case of $\R^2$.

\marginfig{\begin{tikzpicture}[scale=.5]
  \draw[vec] (0,0)--node[below right=-2pt,font=\scriptsize]{$u$} (2.6,.6);
  \draw[vec] (2.6,.6)--node[right,font=\scriptsize]{$v$} (3.3,2.6);
  \draw[vec] (0,0)--node[above left=-2pt,font=\scriptsize]{$u+v$} (3.3,2.6);
\end{tikzpicture}}{The sum of two vectors.}%
Suppose we have two vectors $u$ and $v$ in $\R^2$ that we want to add. Move the
vector $v$ parallel to itself so that its initial point coincides with the end
point of the vector $u$, as shown here. The sum $u+v$ then equals the vector
whose initial point equals the initial point of $u$ and whose end point equals
the end point of the vector $v$, as shown here.

In the next definition, the $0$ on the right side of the displayed equation is
the list $0\in\F^n$.

\begin{definition}[Additive inverse in $\F^n$, $-x$]\label{ax:1.17}
For $x\in\F^n$, the \emph{additive inverse} of $x$, denoted by $-x$, is the
vector $-x\in\F^n$ such that
\[ x+(-x)=0. \]
Thus if $x=(x_1,\dots,x_n)$, then $-x=(-x_1,\dots,-x_n)$.
\end{definition}

\marginfig{\begin{tikzpicture}[scale=.5]
  \draw[vec] (0.4,0.9)--node[above left=-2pt,font=\scriptsize]{$x$} (2.0,2.0);
  \draw[vec] (2.6,1.1)--node[below right=-2pt,font=\scriptsize]{$-x$} (1.0,0);
\end{tikzpicture}}{A vector and its additive inverse.}%
The additive inverse of a vector in $\R^2$ is the vector with the same length
but pointing in the opposite direction. The figure here illustrates this way of
thinking about the additive inverse in $\R^2$. As you can see, the vector
labeled $-x$ has the same length as the vector labeled $x$ but points in the
opposite direction.

Having dealt with addition in $\F^n$, we now turn to multiplication. We could
define a multiplication in $\F^n$ in a similar fashion, starting with two
elements of $\F^n$ and getting another element of $\F^n$ by multiplying
corresponding coordinates. Experience shows that this definition is not useful
for our purposes. Another type of multiplication, called scalar multiplication,
will be central to our subject. Specifically, we need to define what it means
to multiply an element of $\F^n$ by an element of $\F$.

\begin{definition}[Scalar multiplication in $\F^n$]\label{ax:1.18}
The \emph{product} of a number $\lambda$ and a vector in $\F^n$ is computed by
multiplying each coordinate of the vector by $\lambda$:
\[ \lambda(x_1,\dots,x_n) = (\lambda x_1,\dots,\lambda x_n); \]
here $\lambda\in\F$ and $(x_1,\dots,x_n)\in\F^n$.
\end{definition}

\marginnote{Scalar multiplication in $\F^n$ multiplies together a scalar and a
vector, getting a vector. In contrast, the dot product in $\R^2$ or $\R^3$
multiplies together two vectors and gets a scalar. Generalizations of the dot
product will become important in Chapter~6.}%
Scalar multiplication has a nice geometric interpretation in $\R^2$. If
$\lambda>0$ and $x\in\R^2$, then $\lambda x$ is the vector that points in the
same direction as $x$ and whose length is $\lambda$ times the length of $x$. In
other words, to get $\lambda x$, we shrink or stretch $x$ by a factor of
$\lambda$, depending on whether $\lambda<1$ or $\lambda>1$.

\marginfig[2mm]{\begin{tikzpicture}[scale=.5]
  \draw[vec] (0,0)--node[above left=-2pt,font=\scriptsize]{$x$} (1.6,1.2);
  \draw[vec] (1.2,-.6)--node[below right=-2pt,font=\scriptsize]{$(1/2)x$} (2,0);
  \draw[vec] (2.4,1.6)--node[above left=-2pt,font=\scriptsize]{$(3/2)x$} (4.8,3.4);
  \draw[vec] (4.6,.6)--node[below right=-2pt,font=\scriptsize]{$(-1/2)x$} (3.8,0);
\end{tikzpicture}}{Scalar multiplication.}%
If $\lambda<0$ and $x\in\R^2$, then $\lambda x$ is the vector that points in
the direction opposite to that of $x$ and whose length is $\abs{\lambda}$ times
the length of $x$, as shown here.

\subsection{Digression on Fields}

A \emph{field} is a set containing at least two distinct elements called $0$
and $1$, along with operations of addition and multiplication satisfying all
properties listed in \ref{ax:1.3}. Thus $\R$ and $\C$ are fields, as is the set
of rational numbers along with the usual operations of addition and
multiplication. Another example of a field is the set $\{0,1\}$ with the usual
operations of addition and multiplication except that $1+1$ is defined to
equal $0$.

In this book we will not deal with fields other than $\R$ and $\C$. However,
many of the definitions, theorems, and proofs in linear algebra that work for
the fields $\R$ and $\C$ also work without change for arbitrary fields. If you
prefer to do so, throughout much of this book (except for Chapters~6 and~7,
which deal with inner product spaces) you can think of $\F$ as denoting an
arbitrary field instead of $\R$ or $\C$. For results (except in the inner
product chapters) that have as a hypothesis that $\F$ is $\C$, you can probably
replace that hypothesis with the hypothesis that $\F$ is an algebraically
closed field, which means that every nonconstant polynomial with coefficients
in $\F$ has a zero. A few results, such as Exercise~13 in Section~1C, require
the hypothesis on $\F$ that $1+1\ne0$.

% ------------------------------------------------------------------ exercises
\begin{exc}

\begin{exercise}
Show that $\alpha+\beta=\beta+\alpha$ for all $\alpha,\beta\in\C$.
\end{exercise}
\begin{solution}
Let $\alpha=a+bi$ and $\beta=c+di$ for some $a,b,c,d\in\R$. Then
\begin{align*}
\alpha+\beta &= (a+bi)+(c+di)\\
             &= (a+c)+(b+d)i\\
             &= (c+a)+(d+b)i\\
             &= (c+di)+(a+bi) = \beta+\alpha. \qedhere
\end{align*}
\end{solution}

\begin{exercise}
Show that $(\alpha+\beta)+\lambda=\alpha+(\beta+\lambda)$ for all
$\alpha,\beta,\lambda\in\C$.
\end{exercise}
\begin{solution}
Let $\alpha=a+bi$, $\beta=c+di$, and $\lambda=e+fi$ for some
$a,b,c,d,e,f\in\R$. Then
\begin{align*}
(\alpha+\beta)+\lambda &= \bigl((a+bi)+(c+di)\bigr)+(e+fi)\\
  &= \bigl((a+c)+(b+d)i\bigr)+(e+fi)\\
  &= (a+c+e)+(b+d+f)i\\
  &= \bigl(a+(c+e)\bigr)+\bigl(b+(d+f)\bigr)i\\
  &= (a+bi)+\bigl((c+di)+(e+fi)\bigr) = \alpha+(\beta+\lambda). \qedhere
\end{align*}
\end{solution}

\begin{exercise}
Show that $(\alpha\beta)\lambda=\alpha(\beta\lambda)$ for all
$\alpha,\beta,\lambda\in\C$.
\end{exercise}
\begin{solution}
Let $\alpha=a+bi$, $\beta=c+di$, and $\lambda=e+fi$ for some
$a,b,c,d,e,f\in\R$. Then
\begin{align*}
(\alpha\beta)\lambda &= \bigl((a+bi)(c+di)\bigr)(e+fi)\\
  &= \bigl((ac-bd)+(ad+bc)i\bigr)(e+fi)\\
  &= (ace-bde-adf-bcf)+(acf-bdf+ade+bce)i\\
  &= (ace-adf-bde-bcf)+(acf+ade+bce-bdf)i\\
  &= (a+bi)\bigl((ce-df)+(cf+de)i\bigr)\\
  &= (a+bi)\bigl((c+di)(e+fi)\bigr) = \alpha(\beta\lambda). \qedhere
\end{align*}
\end{solution}

\begin{exercise}
Show that $\lambda(\alpha+\beta)=\lambda\alpha+\lambda\beta$ for all
$\lambda,\alpha,\beta\in\C$.
\end{exercise}
\begin{solution}
Let $\alpha=a+bi$, $\beta=c+di$, and $\lambda=e+fi$ for some
$a,b,c,d,e,f\in\R$. Then
\begin{align*}
\lambda(\alpha+\beta) &= (e+fi)\bigl((a+bi)+(c+di)\bigr)\\
  &= (e+fi)\bigl((a+c)+(b+d)i\bigr)\\
  &= \bigl(e(a+c)-f(b+d)\bigr)+\bigl(e(b+d)+f(a+c)\bigr)i\\
  &= (ea-fb+ec-fd)+(eb+fa+ed+fc)i\\
  &= (ea-fb)+(eb+fa)i+(ec-fd)+(ed+fc)i\\
  &= (e+fi)(a+bi)+(e+fi)(c+di) = \lambda\alpha+\lambda\beta. \qedhere
\end{align*}
\end{solution}

\begin{exercise}
Show that for every $\alpha\in\C$, there exists a unique $\beta\in\C$ such that
$\alpha+\beta=0$.
\end{exercise}
\begin{solution}
Let $\alpha=a+bi$ for some $a,b\in\R$. We want to find $\beta=c+di$ such that
$\alpha+\beta=0$. This gives
\[ (a+bi)+(c+di) = (a+c)+(b+d)i = 0+0i. \]
Solving for $c$ and $d$, we get $c=-a$ and $d=-b$. Hence,
$\beta=c+di=-a-bi=-\alpha$.

To see that $\beta$ is unique, suppose there is another $\gamma\in\C$ such that
$\alpha+\gamma=0$. Then
\[ \beta=\beta+0=\beta+(\alpha+\gamma)=(\beta+\alpha)+\gamma=0+\gamma=\gamma. \]
Hence, $\beta$ is unique.
\end{solution}

\begin{exercise}
Show that for every $\alpha\in\C$ with $\alpha\ne0$, there exists a unique
$\beta\in\C$ such that $\alpha\beta=1$.
\end{exercise}
\begin{solution}
Let $\alpha=a+bi$ for some $a,b\in\R$ with $\alpha\ne0$. Then at least one of
$a$ or $b$ is nonzero. We want to find $\beta=c+di$ such that $\alpha\beta=1$.
If $a=0$, then $b\ne0$, and $\beta=1/(bi)=-i/b$. If $b=0$, then $a\ne0$, and
$\beta=1/a$. Otherwise, both $a$ and $b$ are nonzero. We then have
\[ (a+bi)(c+di) = (ac-bd)+(ad+bc)i = 1+0i. \]
Solving for $c$ and $d$, we get the system of equations
\[ \begin{cases} ac-bd=1\\ ad+bc=0 \end{cases}
   \implies
   \begin{cases} a^2c-abd=a\\ b^2c+abd=0. \end{cases} \]
Adding the two equations, we get
\[ (a^2+b^2)c=a \implies c=\frac{a}{a^2+b^2}. \]
Substituting this into the second equation, we get
\[ abd+b^2\frac{a}{a^2+b^2}=0 \implies d=\frac{-b}{a^2+b^2}. \]
Hence,
\[ \beta=c+di=\frac{a}{a^2+b^2}+\frac{-b}{a^2+b^2}\,i. \]
If $\gamma\in\C$ is another complex number such that $\alpha\gamma=1$, then
\[ \gamma=\gamma1=\gamma(\alpha\beta)=(\gamma\alpha)\beta=1\beta=\beta. \]
Hence, $\beta$ is unique.
\end{solution}

\begin{exercise}
Show that
\[ \frac{-1+\sqrt3\,i}{2} \]
is a cube root of $1$ (meaning that its cube equals $1$).
\end{exercise}
\begin{solution}
Note that the complex number given can be written as
\[ 1/2\bigl(-1+\sqrt3\,i\bigr). \]
Cubing $-1+\sqrt3\,i$, we get
\begin{align*}
\bigl(-1+\sqrt3\,i\bigr)^3
  &= (-1)^3+3(-1)^2\bigl(\sqrt3\,i\bigr)+3(-1)\bigl(\sqrt3\,i\bigr)^2+\bigl(\sqrt3\,i\bigr)^3\\
  &= -1+3\sqrt3\,i+3(-1)(-3)+\bigl(\sqrt3\,i\bigr)^3\\
  &= -1+3\sqrt3\,i+9-3\sqrt3\,i\\
  &= 8.
\end{align*}
Therefore,
\[ \biggl(\frac{-1+\sqrt3\,i}{2}\biggr)^{\!3}
   = \frac{\bigl(-1+\sqrt3\,i\bigr)^3}{8} = 8/8 = 1. \qedhere \]
\end{solution}

\begin{exercise}
Find two distinct square roots of $i$.
\end{exercise}
\begin{solution}
Let $z=a+bi$ for some $a,b\in\R$. We want to find $z$ such that $z^2=i$. This
gives
\[ (a+bi)^2 = a^2-b^2+2abi = 0+1i. \]
Solving for $a$ and $b$, we get the system of equations
\[ \begin{cases} a^2-b^2=0\\ 2ab=1. \end{cases} \]
The first equation yields $a=b$ or $a=-b$. If $a=b$, then the second equation
gives
\[ 2b^2=1 \implies b=\pm1/\sqrt2 \implies a=\pm1/\sqrt2. \]
If $a=-b$, this results in $-2b^2=1$, which has no solution when $b\in\R$.
Hence, the two distinct square roots of $i$ are
\[ 1/\sqrt2+1/\sqrt2\,i \quad\text{and}\quad
   -1/\sqrt2-1/\sqrt2\,i. \qedhere \]
\end{solution}

\begin{exercise}
Find $x\in\R^4$ such that
\[ (4,-3,1,7)+2x=(5,9,-6,8). \]
\end{exercise}
\begin{solution}
Let $x=(x_1,x_2,x_3,x_4)\in\R^4$. The equation becomes
\[ (4,-3,1,7)+2(x_1,x_2,x_3,x_4)=(5,9,-6,8), \]
which simplifies to the system of equations
\[ \begin{cases} 5=4+2x_1\\ 9=-3+2x_2\\ -6=1+2x_3\\ 8=7+2x_4 \end{cases}
   \implies
   \begin{cases} 2x_1=1\\ 2x_2=12\\ 2x_3=-7\\ 2x_4=1. \end{cases} \]
Therefore,
\[ x=\Bigl(1/2,\,6,\,-7/2,\,1/2\Bigr). \qedhere \]
\end{solution}

\begin{exercise}
Explain why there does not exist $\lambda\in\C$ such that
\[ \lambda(2-3i,\,5+4i,\,-6+7i)=(12-5i,\,7+22i,\,-32-9i). \]
\end{exercise}
\begin{solution}
Suppose there exists $\lambda\in\C$ such that
\[ \lambda(2-3i,\,5+4i,\,-6+7i)=(12-5i,\,7+22i,\,-32-9i). \]
This gives the system of equations
\[ \begin{cases} \lambda(2-3i)=12-5i\\ \lambda(5+4i)=7+22i\\ \lambda(-6+7i)=-32-9i. \end{cases} \]
Solving the first equation for $\lambda$, we get
\[ \lambda=\frac{12-5i}{2-3i}=\frac{(12-5i)(2+3i)}{(2-3i)(2+3i)}
   =\frac{24+36i-10i-15i^2}{4+9}=\frac{24+26i+15}{13}=3+2i. \]
However, using this value for $\lambda$ in the third equation, we get
\[ (3+2i)(-6+7i)=-18+21i-12i+14i^2=-18+9i-14=-32+9i\ne-32-9i. \]
Hence, no such $\lambda$ exists.
\end{solution}

\begin{exercise}
Show that $(x+y)+z=x+(y+z)$ for all $x,y,z\in\F^n$.
\end{exercise}
\begin{solution}
Let $x=(x_1,\dots,x_n)$, $y=(y_1,\dots,y_n)$, and $z=(z_1,\dots,z_n)$. Then
\begin{align*}
(x+y)+z &= \bigl((x_1,\dots,x_n)+(y_1,\dots,y_n)\bigr)+(z_1,\dots,z_n)\\
  &= (x_1+y_1,\dots,x_n+y_n)+(z_1,\dots,z_n)\\
  &= \bigl((x_1+y_1)+z_1,\dots,(x_n+y_n)+z_n\bigr)\\
  &= \bigl(x_1+(y_1+z_1),\dots,x_n+(y_n+z_n)\bigr)\\
  &= (x_1,\dots,x_n)+\bigl((y_1,\dots,y_n)+(z_1,\dots,z_n)\bigr)\\
  &= x+(y+z). \qedhere
\end{align*}
\end{solution}

\begin{exercise}
Show that $(ab)x=a(bx)$ for all $x\in\F^n$ and all $a,b\in\F$.
\end{exercise}
\begin{solution}
Let $x=(x_1,\dots,x_n)\in\F^n$ and $a,b\in\F$. Then
\begin{align*}
(ab)x &= (ab)(x_1,\dots,x_n)\\
      &= \bigl((ab)x_1,\dots,(ab)x_n\bigr)\\
      &= \bigl(a(bx_1),\dots,a(bx_n)\bigr)\\
      &= a(bx_1,\dots,bx_n)\\
      &= a(bx). \qedhere
\end{align*}
\end{solution}

\begin{exercise}
Show that $1x=x$ for all $x\in\F^n$.
\end{exercise}
\begin{solution}
Let $x=(x_1,\dots,x_n)\in\F^n$. Then
\begin{align*}
1x &= 1(x_1,\dots,x_n)\\
   &= (1x_1,\dots,1x_n)\\
   &= (x_1,\dots,x_n)\\
   &= x. \qedhere
\end{align*}
\end{solution}

\begin{exercise}
Show that $\lambda(x+y)=\lambda x+\lambda y$ for all $\lambda\in\F$ and all
$x,y\in\F^n$.
\end{exercise}
\begin{solution}
Let $x=(x_1,\dots,x_n)$ and $y=(y_1,\dots,y_n)$ be elements of $\F^n$ and
$\lambda\in\F$. Then
\begin{align*}
\lambda(x+y) &= \lambda\bigl((x_1,\dots,x_n)+(y_1,\dots,y_n)\bigr)\\
  &= \lambda(x_1+y_1,\dots,x_n+y_n)\\
  &= \bigl(\lambda(x_1+y_1),\dots,\lambda(x_n+y_n)\bigr)\\
  &= (\lambda x_1+\lambda y_1,\dots,\lambda x_n+\lambda y_n)\\
  &= (\lambda x_1,\dots,\lambda x_n)+(\lambda y_1,\dots,\lambda y_n)\\
  &= \lambda(x_1,\dots,x_n)+\lambda(y_1,\dots,y_n) = \lambda x+\lambda y. \qedhere
\end{align*}
\end{solution}

\begin{exercise}
Show that $(a+b)x=ax+bx$ for all $a,b\in\F$ and all $x\in\F^n$.
\end{exercise}
\begin{solution}
Let $x=(x_1,\dots,x_n)\in\F^n$ and $a,b\in\F$. Then
\begin{align*}
(a+b)x &= (a+b)(x_1,\dots,x_n)\\
       &= \bigl((a+b)x_1,\dots,(a+b)x_n\bigr)\\
       &= (ax_1+bx_1,\dots,ax_n+bx_n)\\
       &= (ax_1,\dots,ax_n)+(bx_1,\dots,bx_n)\\
       &= ax+bx. \qedhere
\end{align*}
\end{solution}
\end{exc}

\axquote{``Can you do addition?'' the White Queen asked. ``What's one and one
and one and one and one and one and one and one and one and one?''\\
``I don't know,'' said Alice. ``I lost count.''}{\emph{Through the Looking
Glass}, Lewis Carroll}
